% from ega4-2-fr.tex
\oldpage[IV-2]{230}

\EGAIVtwoContentsLine{4.7.}{Multiplicités dans la décomposition primaire sur un préschéma algébrique}{75}
\EGAIVtwoContentsLine{4.8.}{Corps de définition}{80}
\EGAIVtwoContentsLine{4.9.}{Corps de définition d'une partie d'un préschéma}{84}

\EGAIVtwoContentsLine{\S~5.}{Dimension, profondeur, régularité dans les préschémas localement noethériens}{86}
\EGAIVtwoContentsLine{5.1.}{Dimension des préschémas}{86}
\EGAIVtwoContentsLine{5.2.}{Dimension d'un préschéma algébrique}{90}
\EGAIVtwoContentsLine{5.3.}{Dimension du support d'un Module et polynôme de Hilbert}{92}
\EGAIVtwoContentsLine{5.4.}{Dimension de l'image d'un morphisme}{93}
\EGAIVtwoContentsLine{5.5.}{Formule des dimensions pour un morphisme de type fini}{94}
\EGAIVtwoContentsLine{5.6.}{Formule des dimensions et anneaux universellement caténaires}{97}
\EGAIVtwoContentsLine{5.7.}{Profondeur et propriété $(S_k)$}{103}
\EGAIVtwoContentsLine{5.8.}{Préschémas réguliers et propriété $(R_k)$. Critère de normalité de Serre}{107}
\EGAIVtwoContentsLine{5.9.}{Modules $Z$-purs et $Z$-clos}{109}
\EGAIVtwoContentsLine{5.10.}{Propriété $(S_2)$ et $Z$-clôture}{114}
\EGAIVtwoContentsLine{5.11.}{Critères de cohérence pour les Modules $\mathcal H^0_{X/Z}(\mathscr F)$}{122}
\EGAIVtwoContentsLine{5.12.}{Relations entre les propriétés d'un anneau local noethérien $A$ et d'un anneau quotient $A/tA$}{126}
\EGAIVtwoContentsLine{5.13.}{Propriétés de permanence par passage à la limite inductive}{131}

\EGAIVtwoContentsLine{\S~6.}{Morphismes plats de préschémas localement noethériens}{134}
\EGAIVtwoContentsLine{6.1.}{Platitude et dimension}{135}
\EGAIVtwoContentsLine{6.2.}{Platitude et dimension projective}{137}
\EGAIVtwoContentsLine{6.3.}{Platitude et profondeur}{138}
\EGAIVtwoContentsLine{6.4.}{Platitude et propriété $(S_k)$}{141}
\EGAIVtwoContentsLine{6.5.}{Platitude et propriété $(R_k)$}{143}
\EGAIVtwoContentsLine{6.6.}{Propriétés de transitivité}{145}
\EGAIVtwoContentsLine{6.7.}{Application aux changements de base dans les préschémas algébriques}{145}
\EGAIVtwoContentsLine{6.8.}{Morphismes réguliers, normaux, réduits, lisses}{150}
\EGAIVtwoContentsLine{6.9.}{Le théorème de platitude générique}{153}
\EGAIVtwoContentsLine{6.10.}{Dimension et profondeur d'un Module normalement plat le long d'un sous-préschéma fermé}{155}
\EGAIVtwoContentsLine{6.11.}{Critères pour que les ensembles $U_{S_n}(\mathscr F)$ ou $U_{C_n}(\mathscr F)$ soient ouverts}{158}
\EGAIVtwoContentsLine{6.12.}{Critères de Nagata pour que $\operatorname{Reg}(X)$ soit ouvert}{163}
\EGAIVtwoContentsLine{6.13.}{Critères pour que $\operatorname{Nor}(X)$ soit ouvert}{168}
\EGAIVtwoContentsLine{6.14.}{Changement de base et clôture intégrale}{169}
\EGAIVtwoContentsLine{6.15.}{Préschémas géométriquement unibranches}{176}

\clearpage
